%; whizzy paragraph -pdf xpdf -latex ./whizzypdfptex.sh
%; whizzy-paragraph "^\\\\begin{frame}"
% latex beamer presentation.
% platex, latex-beamer でコンパイルすることを想定。 

%     Tokyo Debian Meeting resources
%     Copyright (C) 2009 Junichi Uekawa

%     This program is free software; you can redistribute it and/or modify
%     it under the terms of the GNU General Public License as published by
%     the Free Software Foundation; either version 2 of the License, or
%     (at your option) any later version.

%     This program is distributed in the hope that it will be useful,
%     but WITHOUT ANY WARRANTY; without even the implied warranty of
%     MERCHANTABILITY or FITNESS FOR A PARTICULAR PURPOSE.  See the
%     GNU General Public License for more details.

%     You should have received a copy of the GNU General Public License
%     along with this program; if not, write to the Free Software
%     Foundation, Inc., 51 Franklin St, Fifth Floor, Boston, MA  02110-1301 USA

\documentclass[cjk,dvipdfmx,12pt]{beamer}
\usetheme{Tokyo}
\usepackage{monthlypresentation}

%  preview (shell-command (concat "evince " (replace-regexp-in-string "tex$" "pdf"(buffer-file-name)) "&"))
%  presentation (shell-command (concat "xpdf -fullscreen " (replace-regexp-in-string "tex$" "pdf"(buffer-file-name)) "&"))
%  presentation (shell-command (concat "evince " (replace-regexp-in-string "tex$" "pdf"(buffer-file-name)) "&"))

%http://www.naney.org/diki/dk/hyperref.html
%日本語EUC系環境の時
\AtBeginDvi{\special{pdf:tounicode EUC-UCS2}}
%シフトJIS系環境の時
%\AtBeginDvi{\special{pdf:tounicode 90ms-RKSJ-UCS2}}

\title{東京エリア Debian 勉強会}
\subtitle{資料}
\author{上川 純一 dancer@debian.or.jp\\IRC nick: dancerj}
\date{2009年3月21日}
\logo{\includegraphics[width=8cm]{image200607/openlogo-light.eps}}

\begin{document}

\frame{\titlepage{}}

\emtext{設営準備にご協力ください}

\begin{frame}{今日の参加の目標 1/3}

\begin{itemize}
 \item 岩松: Debianで尊敬する日々野さんの lisp について勉強したい
 \item あけど: 日本のlispの権威である日々野さんから　Debian で lisp を教
       わりたい
 \item 前田:最近　erlang を勉強している、ところで日々野さんの common lisp について
       語るというので聞きたい。はやく日本 ocaml ユーザ会を立ち上げてくだ
       さい。
 \item キタハラ: 最近 on lisp という本があたったので流行っているのにちがいな
       い。
 \item たかはし: Debianの人たちがどうやってつかってるのかな? が知りたい。
 \item キタムラ: 日比野さんをリスペクトして参加。
\end{itemize}
\end{frame}

\begin{frame}{今日の参加の目標 2/3}
\begin{itemize}
 \item じつかた: プライベートなところではDebianよ。
 \item やまだ:　やっと Debian で Git　使い始めたので、日々野さんのlisp魂
       の薫陶を(以下略
 \item 日比野: lisp に興味をもってください。
 \item 藤澤とおる: Debianでパッケージを作るという話をするけど、つっこみ大好き。
 \item よしの: 普段とおいところなので参加できてなかったが、近いので今日
       参加しました。
 \item 虎: サーバにDebianつかってます。新しいリリースはすばらしい。
       haskell と ocamlつかってます。 lispに興味あります。
 \item こたに: 会社と研究室でたまに Debian つかってます。
 \item まつう: gentoo つかってます。仕事ではたまにDebianです。package 作
       成に興味あります。
 \item かい:　プライベート隠れDebianです。開発者のお話がききたい。今日は日比野
       さんにたぶらかされて参加。
\end{itemize}
\end{frame}

\begin{frame}{今日の参加の目標 3/3}
 \begin{itemize}
  \item たなか: 新潟からきました。
  \item ささき: uwabamiです。関西からきました。雰囲気を見学したい。
  \item まとはら: koedolugからきました。2回目。 lisp やったことないので、
	lisp を覚えたい。
  \item こう: Debian の人にあいたいとおもってきた。
  \item John: 日本語わかりません。 lisp 使ってます。Debianも長い間つかっ
	ています。
  \item よしだ@板橋: lisp を勉強にきたふりをして安田講堂を見物。
  \item 藤澤りそう: 参加したのは習慣になっているので、いつもの癖です。目標はいまから考えます。
 \end{itemize}
\end{frame}

\section{}
\begin{frame}
 \frametitle{Agenda}
\begin{minipage}[t]{0.45\hsize}
  \begin{itemize}
  \item 注意事項
	\begin{itemize}
	 \item ?
	\end{itemize}
  \item 最近あったDebian関連のイベント報告
	\begin{itemize}
	 \item 前回の勉強会
	 \item OSC
	 \item Ubuntu オフラインミーティング＠秋葉原
	 \item Linux Consortium 10 years event
	 \item Hack Cafe
	\end{itemize}
 \end{itemize}
\end{minipage} 
\begin{minipage}[t]{0.45\hsize}
 \begin{itemize}
  \item Debian クイズ
  \item 研究室のソフトウェアを Debian パッケージにしてみる
  \item Debian での Common Lispプログラミング環境
 \end{itemize}
\end{minipage}
\end{frame}

\section{最近}

\begin{frame}
 \frametitle{2009年02月}
\begin{minipage}[t]{0.45\hsize}
  \begin{itemize}
  \item 注意事項
	\begin{itemize}
	 \item 飲食禁止
	 \item 政治/宗教/営利活動禁止
	\end{itemize}
  \item 最近のDebian関連のイベント
	\begin{itemize}
	 \item Lenny release
	 \item Linux Consortium 10 Years Event !!
	 \item Debian ハックカフェ開始
	 \item 第22回関西エリアDebian勉強会
	\end{itemize}
 \end{itemize}
\end{minipage} 
\begin{minipage}[t]{0.45\hsize}
	\begin{itemize}
	 \item OSC2009 Tokyo Spring
	 \item Debianパッケージングハンズオン
	 \item Debian マシン展示(eeepc/chamby)
	 \item グッズ販売
	\end{itemize}
\end{minipage}
\end{frame}

\begin{frame}{Linux Consortium 10 years event}
 
\end{frame}

\begin{frame}{Ubuntu オフラインミーティング＠秋葉原}
 
\end{frame}

\begin{frame}{Hack Cafe}

毎週水曜日、週に一回東京のどっかのカフェでハック。

\end{frame}

\begin{frame}{関西 Debian 勉強会}

岩松さんが同志社大学にて Git についてディスカッション。

\end{frame}

\section{DWN quiz}
\begin{frame}{Debian 常識クイズ}

Debian の常識、もちろん知ってますよね?
知らないなんて恥ずかしくて、知らないとは言えないあんなことやこんなこと、
みんなで確認してみましょう。

今回の出題範囲は\url{debian-devel-announce@lists.debian.org} に投稿された
内容とDebian Project Newsからです。

\end{frame}

\subsection{問題}
%; whizzy-master ../debianmeetingresume200903.tex
% 以上の設定をしているため、このファイルで M-x whizzytex すると、whizzytexが利用できます。

 \santaku
 {2月14日にリリースされたのは}
 {etch}
 {sarge}
 {lenny}
 {C}
{}

 \santaku
 {3月12日にリリースされた Debian Policy のバージョンは}
 {3.8.0}
 {3.8.1}
 {3.8.2}
 {B}
{}

 \santaku
 {Debian Data Export は何か}
 {Debianパッケージに関するデータを容易に取り寄せるためのインターフェイス}
 {全世界にある、あらゆるデータをDebianパッケージ化するプロジェクト}
 {DebianのあらゆるデータをGoogleにすべて喰わせてみましたというプロジェクト}
 {A}
{}

 \santaku
 {Olivier BergerがMLに流したbts-linkに関する情報は}
 {bts-linkのソース管理リポジトリが壊れました}
 {bts-linkでリンクできるBTSを増やしました}
 {bts-linkサービスは終了です}
 {B}
{}

 \santaku
 {Debian FTP masterが変わりました。誰が新しく入ったでしょうか。}
 {Kei Hibino}
 {Ryan Murray}% FTP master を辞めた人
 {Mark Hymers}% 新しく入った人
 {C}
{}

 \santaku
 {3月16日にJoerg JaspertがMLに流したアナウンスは}
 {Debian の新しいロゴが決まりました}
 {パッケージのセクションが追加されます}
 {いくつかのディストリビューションと吸収合併します}
 {B}
{}

 \santaku
 {3月21日にリリースされた pbuilder のバージョンは?}
 {0.187}
 {1.298}
 {3.1}
 {A}
{}

 \santaku
 {Debian.org DPL に立候補していないのは}
 {Stefano Zacchiroli}
 {Steve McIntyre}
 {Nobuhiro Iwamatsu}
 {C}
{}

 \santaku
 {Debian JP 選挙の立候補しめきりは?}
 {3月20日}
 {3月22日}
 {3月24日}
 {B}
{}


\begin{frame}{2009年計画}

{\scriptsize
 \begin{enumerate}
  \item 新年の企画 (アンサンブル荻窪開催)
  \item OSC Tokyo
  \item VAIO P インストール記録、
	カーネル読書会　ディストリビューション大集合(小林さん)(東京大学?)
  \item Git Handson (岩松)(あんさんぶる荻窪?)
  \item 家Debianサーバ vs 職場のネットワーク(千代田区都立図書館?\footnote{\url{http://www.library.chiyoda.tokyo.jp/}})
  \item Asterisk (東京大学?)
  \item スペインにて開催
  \item Debconf報告会
  \item OSC Fall?
  \item udev + HAL
  \item 3D graphics 開発 
  \item Debian サーバ+VMware + 各種OS、
	他の仮想化ツール(vserver etc.)、
	忘年会
 \end{enumerate}
}
\end{frame}

\begin{frame}{2009年からのフロー(案)}
\begin{itemize}
 \item 前回の勉強会で発表内容を打ち合わせ　(2/21)
 \item 二週間後にドラフト作成 (3/7)
 \item 三週間後に資料のテクニカルレビュー・メーリングリストで議論 (3/8-3/14)
 \item 印刷発注準備 (3/15-3/19)
 \item 四週間後に発表 (3/21)
\end{itemize}
\end{frame}

\emtext{研究室のソフトウェアを Debian パッケージにしてみる}

\emtext{Debian での Common Lispプログラミング環境}

\emtext{次回の勉強会}
\begin{frame}{次回の勉強会}
次回 は 2008年4月18日に行われます。
荻窪のあんさんぶる荻窪で開催予定です。

ネタ募集中。

\begin{itemize}
 \item 事前課題: 「やまねさんの開発効率を上げる方法を提案してください。」
 \item 藤澤さん:　連載GXP: ITPしてみました。
 \item John: I made a clozurecl package!
 \item やまね: 俺の開発効率をあげてください session。
\end{itemize}
\end{frame}

\begin{frame}{Debconf 2009}
\begin{itemize}
 \item スペイン参加
 \item 飛行機代20万弱
 \item 7/23-7/30
\end{itemize}
\end{frame}

\begin{frame}{宴会場所}
\begin{itemize}
 \item 宴会場所\\
       本日の宴会は「鉄板焼き 郷」です。\\
       18:30 開始です。
       参加者は全員で移動しましょう。
 \item 片付け\\
       部屋を片付けるのにご協力ください。
\end{itemize}
\end{frame}

\end{document}

;;; Local Variables: ***
;;; outline-regexp: "\\([ 	]*\\\\\\(documentstyle\\|documentclass\\|emtext\\|section\\|begin{frame}\\)\\*?[ 	]*[[{]\\|[]+\\)" ***
;;; End: ***
